\documentclass[10pt,a4paper]{amsart}
\usepackage[margin=19mm]{geometry}
\usepackage{amsmath,amssymb,mathtools}
\usepackage[scr=rsfs,cal=euler]{mathalfa}
\usepackage{xfrac}
\usepackage[unicode,hidelinks,bookmarks=false]{hyperref}
\newcommand{\ie}{i.e.\ }
\newcommand{\cf}{cf.\ }
\newcommand{\resp}{respectively\ }
\newcommand{\Subsection}{\subsection}
\providecommand{\nobreakdash}{\nobreak\hbox{-}}
\setlength{\emergencystretch}{2em}
\multlinegap0pt
\usepackage{amssymb}
\usepackage[shortlabels]{enumitem}
\usepackage{mathtools}
\usepackage{tikz}
\usepackage{xcolor}
\newtheorem{lemma}{Lemma}
\newtheorem{thm}{Theorem}
\title{On numerical semigroups with embedding dimension four}
\author{Kazimierz Chomicz}
\date{Source: arXiv:2604.25653v3 (CC BY 4.0)}
\begin{document}
\maketitle

\noindent\emph{Source and licence.} On numerical semigroups with embedding dimension four. Version arXiv:2604.25653v3 (CC BY 4.0), distributed by arXiv under the Creative Commons Attribution 4.0 International licence, http://creativecommons.org/licenses/by/4.0/, as recorded in the arXiv licence metadata for this exact version and shown on the abstract page https://arxiv.org/abs/2604.25653v3. Full licence notice: \texttt{licenses/arxiv-2604.25653-CC-BY-4.0.txt}.\par\smallskip

\noindent\textit{Editorial scope.} Counted unit: the article's thm theoremprocedure (source0.tex lines 265--274). The article's complete original proof of it is delivered with it (source0.tex lines 275--316, one \texttt{proof} environment of 42 lines). No mathematics is written, repaired or re-derived: the statement and the proof are the article's own lines, in the article's own notation, and no retained line is substituted or re-flowed.  Retained uncounted as the source-local context the proof needs: lines 226--231; lines 241--243; lines 252--254.  Nothing was omitted from any retained span, so the package carries no omission record.  No retained line contains a citation, so the wrapper carries no bibliography and no bibliography entry.  No retained line contains a figure environment, an image-inclusion command or drawing code, so no image is delivered and the package declares no asset dependency.  Editorial formatting only, in addition: an AMS \texttt{amsart} preamble that restates the article's own declarations and macros used by the retained lines, namely the environments \texttt{lemma}, \texttt{thm} on separate counters, together with the standard packages those lines need.  The retained environments print in this document as Lemma 1, Theorem 1 in order of appearance, whereas the article prints the same items with the numbering of its own full document.  The complete native source of the article as distributed in the arXiv e-print for this exact version is bundled unchanged as \texttt{sources/199318/source0.tex}.
\begin{align}
    a_{00} d_0 &= a_{01} d_1 + a_{02}d_2 + a_{03} d_3,  \label{minrelation-0} \\
    a_{11} d_1 &= a_{10} d_0+a_{12} d_2 + a_{13} d_3,  \\
    a_{22} d_2 &=a_{20} d_0+ a_{21} d_1 + a_{23} d_3,  \\
    a_{33} d_3 &= a_{30} d_0 +a_{31} d_1 + a_{32} d_2, 
\end{align}

\begin{equation}\label{positivepoints}
 P_1=(b_{11},-b_{12},-b_{13}), \ P_2=(-b_{21},b_{22},-b_{23}), \ P_3 = (-b_{31}, -b_{32}, b_{33}),
\end{equation} 

\begin{lemma}\label{1remlemma}
    Deleting from the initial collection of cubes the region associated to a point $P = (\mu_1,\mu_2,\mu_3) \in \mathbb{Z}^3$, satisfying $\mu_1d_1 + \mu_2d_2 + \mu_3d_3 = \mu d_0$ for some positive integer $\mu$, does not remove any cube labeled with an element of the Ap{\'e}ry set with respect to $d_0$. 
\end{lemma}

\begin{thm}\label{theoremprocedure}
     Delete from the initial collection of cubes the regions associated to the points \eqref{positivepoints}, and for every triple $(a_{01},a_{02},a_{03}) \in \mathbb{N}^3$ satisfying \eqref{minrelation-0} delete the region associated to $(a_{01},a_{02},a_{03})$. Additionally, for all identities of the form
    $$ \lambda_0d_0 + \lambda_1d_1=\lambda_2d_2 + \lambda_3d_3,  $$
    with $\lambda_i \in \mathbb{Z}_{>0}$, $\lambda_2 < b_{22}$, and $\lambda_3 < b_{33}$, remove the region associated to the point $(-\lambda_1,\lambda_2,\lambda_3)$. For all identities of the form
    $$ \lambda_0d_0 + \lambda_2d_2=\lambda_1d_1 + \lambda_3d_3,  $$
    with $\lambda_i \in \mathbb{Z}_{>0}$, $\lambda_1 < b_{11}$, and $\lambda_3 < b_{33}$, remove the region associated to the point $(\lambda_1,-\lambda_2,\lambda_3)$. Finally, for all identities of the form
    $$ \lambda_0d_0 + \lambda_3d_3=\lambda_1d_1 + \lambda_2d_2,  $$
    with $\lambda_i \in \mathbb{Z}_{>0}$, $\lambda_1 < b_{11}$, and $\lambda_2 < b_{22}$, remove the region associated to the point $(\lambda_1,\lambda_2,-\lambda_3)$.
    Call the remaining collection of cubes the figure $R$. Then, the set of all cubes from the initial collection of cubes which are labeled with an element of the Ap{\'e}ry set is the figure $R$.
\end{thm}

\begin{proof}
The condition $\lambda_i < b_{ii}$ ensures that we consider only the identities that matter --- the points \eqref{positivepoints} already delete the regions $\{x > b_{11}\}$, $\{y > b_{22}\}$, $\{z > b_{33}\}$. Lemma \ref{1remlemma} ensures that every cube labeled with an element of the Ap{\'e}ry set is in $R$. We need to prove that only the cubes labeled with elements of the Ap{\'e}ry set are in $R$. Suppose, to the contrary, that for some remainder $r \in \{0,1,\dots,d_0-1\}$, $\lambda d_0 + r$ and $\lambda' d_0 + r$, with $\lambda' > \lambda$, are labels in $R$. Now for some $\lambda_1, \lambda_2, \lambda_3,\lambda_1', \lambda_2', \lambda_3' \in \mathbb{N}$ we have
    \begin{equation}\label{eq1theorem}
        \lambda d_0 + r = \lambda_1d_1 + \lambda_2d_2 + \lambda_3d_3,
    \end{equation}
     \begin{equation}\label{eq2theorem}
        \lambda' d_0 + r = \lambda_1'd_1 + \lambda_2'd_2 + \lambda_3'd_3,
    \end{equation}
    with $[[\lambda_1,\lambda_2,\lambda_3]], [[\lambda_1',\lambda_2',\lambda_3']] \in R$.
    Subtracting \eqref{eq1theorem} from \eqref{eq2theorem} we get
    \begin{equation}\label{eq3theorem}
        (\lambda' - \lambda)d_0 = (\lambda_1' - \lambda_1)d_1 + (\lambda_2' - \lambda_2)d_2 + (\lambda_3' - \lambda_3)d_3,
    \end{equation}
    where $(\lambda' - \lambda)>0$, by assumption. Consider cases of how many coefficients on the right-hand side of \eqref{eq3theorem} are positive. The left-hand side is positive, hence, on the right-hand side, at least one of the coefficients is positive. 
    
If exactly one coefficient is positive, assume without loss of generality that $(\lambda_1' - \lambda_1)$ is the one. Moving the non positive integers to the other side gives
    \begin{equation}\label{eq4theorem}
        (\lambda' - \lambda)d_0  + (\lambda_2 - \lambda_2')d_2 + (\lambda_3 - \lambda_3')d_3= (\lambda_1' - \lambda_1)d_1.
    \end{equation} 
Since all the above coefficients are nonnegative and $(\lambda_1' - \lambda_1)$ and $(\lambda' - \lambda) $ are positive, by the minimality of $b_{11}$, we get $b_{11} \leq (\lambda_1' - \lambda_1) \leq \lambda_1'$. This is a contradiction with $[[\lambda_1',\lambda_2',\lambda_3']] \in R$.

If exactly two coefficients on the right-hand side of \eqref{eq3theorem} are positive, without loss of generality, suppose that $(\lambda_1' - \lambda_1)$ and $(\lambda_2' - \lambda_2)$ are the two. Moving $(\lambda_3' - \lambda_3)$ to the other side gets
    \begin{equation}\label{eq5theorem}
        (\lambda' - \lambda)d_0  + (\lambda_3 - \lambda_3')d_3= (\lambda_1' - \lambda_1)d_1  + (\lambda_2' - \lambda_2)d_2.
    \end{equation}
The above identity satisfies the conditions of the current theorem, thus the region $\{x > (\lambda_1' - \lambda_1), y > (\lambda_2' - \lambda_2)\}$ was deleted. Again, this contradicts $[[\lambda_1',\lambda_2',\lambda_3']] \in R$, as $\lambda'_1 \geq (\lambda_1' - \lambda_1)$ and $\lambda'_2 \geq (\lambda_2' - \lambda_2)$.

Now, only the case where all the coefficients in \eqref{eq3theorem} are positive is left. Let us substitute $\mu =(\lambda' - \lambda)  > 0$ and $\mu_i =(\lambda_i' - \lambda_i)  >0$ for $i = 1,2,3$. We have
    \begin{equation}\label{eq6theorem}
         \mu d_0 = \mu_1d_1 + \mu_2d_2 + \mu_3d_3,
    \end{equation}
and by the process of deleting the regions, the cube $[[\mu_1,\mu_2,\mu_3]]$ is also in $R$, as $\mu_i \leq \lambda'_i$ for $i = 1,2,3$. 
From the minimality of $a_{00}$, we get $a_{00} \leq \mu$. Now we subtract $a_{00}d_0 = a_{01}d_1 + a_{02}d_2 + a_{03}d_3$ from \eqref{eq6theorem} to get
\begin{equation}\label{eq7theorem}
    (\mu - a_{00}) d_0 = (\mu_1 - a_{01})d_1 +  (\mu_2 - a_{02})d_2 + (\mu_3 - a_{03})d_3.
\end{equation}
    
If $(\mu - a_{00}) = 0$, then 
the region associated to $(\mu_1,\mu_2,\mu_3)$ had been deleted, which contradicts $[[\mu_1,\mu_2,\mu_3]] \in R$. 

Finally, when $(\mu - a_{00}) > 0$ in \eqref{eq7theorem}, we repeat the same case by case consideration as we did with \eqref{eq3theorem}, which results in all the coefficients in \eqref{eq7theorem} being positive. However, this implies that the cube $[[\mu_1,\mu_2,\mu_3]]$, has been deleted by the region associated to $(a_{01}, a_{02}, a_{03})$, which is a contradiction. 
\end{proof}

\end{document}
